\documentclass{exam-zh}

% Visual regression test for GitHub issues #50, #51, and #52.
% Compile twice from testfiles/ with the workspace package taking precedence:
%   TEXINPUTS=..: latexmk -xelatex -interaction=nonstopmode \
%     -file-line-error -synctex=1 fillin-auto-wrap-regression.tex
% The second invocation should report "Nothing to do".

\examsetup{
  page = {
    size = a3paper,
    show-columnline = true,
  },
  fillin = {
    width = auto,
    width-extra = 1em,
    no-answer-type = none,
    show-answer = false,
    type = line,
  },
}

\AtEndPreamble{\geometry{showframe}}

\definecolor{RegressionGreen}{HTML}{176B45}
\definecolor{RegressionGray}{HTML}{F1F3F5}

\newcommand\regressionnote[1]{%
  \par\smallskip
  \noindent
  \colorbox{RegressionGray}{%
    \parbox{\dimexpr\linewidth-2\fboxsep\relax}{#1}%
  }%
  \par\smallskip
}

\newcommand\oldbehavior[1]{%
  \par\noindent\textbf{原错误现象：}#1\par
}

\newcommand\fixedbehavior[1]{%
  \par\noindent\textcolor{RegressionGreen}{\textbf{修复后：}}#1\par
}

\begin{document}
\raggedbottom

\section*{\texttt{\string\fillin} 自动宽度与换行回归测试}

本文档同时覆盖 GitHub Issue \#50、\#51 和 \#52。页面边框与中间分栏线是验收基准：
所有下划线和标点均应位于所属栏的边框内。本文只运行当前修复后的实现；旧效果
用文字说明，避免为了截图而重新引入有缺陷的代码。

\subsection*{1. 短填空与后续标点作为整体换行（Issue \#52）}

\oldbehavior{短填空与紧随的句号虽然被放入同一个不可拆分盒子，但当行末剩余
空间不足时，TeX 可能因为普通断点的 badness 过高而拒绝换行，导致“下划线
加句号”整个单元越出右侧版心。}

\fixedbehavior{排版前先比较“下划线宽度 + 标点宽度”和当前行剩余宽度；放得下
就保持行内，放不下就让整个单元移到下一行，且不会把标点单独留在行首。}

\regressionnote{下面两行是实际输出，不是示意图。第一行预留约 \texttt{9em}，
应当容纳完整单元；第二行只预留约 \texttt{7em}，完整单元应当一起移到下一行。}

\noindent
\makebox[\dimexpr\linewidth-9em\relax][l]
  {情形 A（仍能放下）：当前行末保留足够宽度}%
\fillin*[生物化学反应]。

\medskip

\noindent
\makebox[\dimexpr\linewidth-7em\relax][l]
  {情形 B（已经放不下）：当前行末剩余宽度不足}%
\fillin*[生物化学反应]。

\regressionnote{修复后的预期效果：情形 A 的下划线和句号留在第一行；情形 B
的下划线和句号同时出现在下一行。两种情形都不越过右侧边框。}

\subsection*{2. 常规自然宽度与标点吸附}

\oldbehavior{早期实现使用固定宽度，或在换行时把逗号、句号留到下一行开头；
段首缩进还可能被重复计入可用宽度。}

\fixedbehavior{\texttt{width=auto} 按隐藏答案的自然宽度留空，另加
\texttt{width-extra}；常见右标点与短填空保持在一起，段首和显式
\texttt{\string\noindent} 分别按各自的真实剩余宽度计算。}

\begin{question}
  组成细胞的化学元素和无机自然界中的化学元素相比，种类
  \fillin*[大体相同]，含量\fillin*[差异很大]。
\end{question}

\begin{question}
  水在细胞中以\fillin*[自由水和结合水]两种形式存在。
\end{question}

\regressionnote{修复后的预期效果：两个填空长度随原答案长度变化；逗号、句号
和短填空保持相邻，没有孤立在下一行开头。}

% In A3 mode, \newpage ends the left column and starts the right column.
\newpage

\section*{右栏：A3 双栏剩余宽度（Issue \#50）}

本节被有意放在 A3 页面的右栏，以运行 Issue \#50 的真实失败路径。

\oldbehavior{\texttt{linegoal} 以左栏边界为基准返回位置；旧计算没有正确补偿
右栏偏移，因而会得到负数或错误的剩余宽度，使自动宽度填空过早换行、长度异常，
甚至越出右栏。}

\fixedbehavior{仅在 A3 双栏的右栏补回栏宽与栏间距；A4 单栏和 A3 左栏的计算
保持不变。下面四题直接采用原问题中的题目。}

\begin{question}
  组成细胞的化学元素和无机自然界中的化学元素相比，种类
  \fillin*[大体相同]，含量\fillin*[差异很大]。
\end{question}

\begin{question}
  组成细胞的大量元素有
  \fillin*[C、H、O、N、P、S、K、Ca、Mg]等，微量元素有
  \fillin*[Fe、Mn、B、Zn、Mo、Cu]等。
\end{question}

\begin{question}
  人体细胞鲜重含量最多的元素是\fillin*[O]，干重含量最多的元素是
  \fillin*[C]。
\end{question}

\begin{question}
  组成细胞元素大多以\fillin*[化合物]的形式存在，组成细胞的无机盐大多以
  \fillin*[离子]的形式存在。
\end{question}

\regressionnote{修复后的预期效果：所有下划线都按答案自然宽度排版；需要换行时
只在右栏内部换行，标点不会落在下一行开头，也不会越过右侧页面边框。}

\subsection*{3. 多轮编译收敛（Issue \#51）}

\oldbehavior{XeTeX 保存的位置有时会在相邻编译轮次间变化 \texttt{1sp}，
\texttt{linegoal} 将这个差异写入辅助文件后，\texttt{latexmk} 可能持续提示
标签变化并反复编译。}

\fixedbehavior{用于剩余宽度计算的位置统一舍入到 \texttt{64sp}，且辅助文件不再
保存未参与宽度计算的原始横坐标。第一次编译完成交叉定位后，第二次执行
\texttt{latexmk} 应报告 \texttt{Nothing to do}。}

\regressionnote{最终验收：日志中没有 \texttt{Overfull \string\hbox}、
没有 \texttt{Package linegoal Warning}；再次执行相同编译命令时无需重编译。}

\end{document}
